% segment-id: o014.aljabr2.chapter9.coendomorphism-coalgebra-b.p001
Untuk funktor $\omega: \mathcal{A} \to \catesubd{Mod}{\mathrm{pfg}}B$, homomorfisme evaluasi $\mathrm{ev}_{\omega(X)}: \omega(X)^\vee \dotimes{\Bbbk} \omega(X) \to B$ merupakan homomorfisme bimodul-$(B, B)$. Homomorfisme kanonik yang ditentukannya tepat merupakan homomorfisme yang didefinisikan sebelumnya,
% segment-id: o014.aljabr2.chapter9.coendomorphism-coalgebra-b.eq001
\begin{equation}\label{eqn:L-counit}
	\begin{aligned}
		\epsilon: \coEnd(\omega) & \to B \\
		[t] & \mapsto \mathrm{ev}_{\omega(X)}(t).
	\end{aligned}
\end{equation}

% segment-id: o014.aljabr2.chapter9.coendomorphism-coalgebra-b.p002
Jika $\omega_1(X)$ bebas, ambil suatu basis-$B$ $(v_i)_{i=1}^n$ dan basis dualnya $(\check{v}_i)_{i=1}^n$. Maka $\lambda_X: \omega_2(X) \to \omega_1(X) \dotimes{B} \coHom(\omega_1, \omega_2)$ dapat ditulis sebagai
% segment-id: o014.aljabr2.chapter9.coendomorphism-coalgebra-b.eq002
\begin{equation}\label{eqn:L-coaction}\begin{aligned}
	w & \mapsto \sum_{i=1}^n v_i \otimes [\check{v}_i \otimes w];
\end{aligned}\end{equation}
% segment-id: o014.aljabr2.chapter9.coendomorphism-coalgebra-b.p003
Berdasarkan hal ini, untuk sembarang $\omega_1, \omega_2, \omega_3$, apabila $\omega_1(X)$ dan $\omega_2(X)$ bebas, $\lambda \in \omega_1(X)^\vee$, dan $w \in \omega_3(X)$, diperoleh
% segment-id: o014.aljabr2.chapter9.coendomorphism-coalgebra-b.eq003
\begin{equation}\label{eqn:L-comult}\begin{aligned}
	\coHom(\omega_1, \omega_3) & \to \coHom(\omega_1, \omega_2) \dotimes{B} \coHom(\omega_2, \omega_3) \\
	[\lambda \otimes w] & \mapsto \sum_i [\lambda \otimes v_i] \otimes [\check{v}_i \otimes w].
\end{aligned}\end{equation}

% segment-id: o014.aljabr2.chapter9.coendomorphism-coalgebra-b.p004
Langkah berikutnya ialah memasukkan struktur monoidal. Untuk itu, $B$ harus komutatif dan kita harus mengambil $\Bbbk = B$ dalam konstruksi sebelumnya. Dengan demikian, $\cated{Mod}B$ dan $\catesubd{Mod}{\mathrm{pfg}}B$ keduanya merupakan kategori monoidal simetris terhadap $\otimes_B$, dengan $B$ sebagai unit. Untuk menegaskan keadaan ini, $\coHom$ dan $\coEnd$ yang bersesuaian akan dinotasikan dengan $\coHom_B$ dan $\coEnd_B$; dalam hal ini, bimodul-$(B, B)$ sama saja dengan modul-$B$.

% segment-id: o014.aljabr2.chapter9.coendomorphism-coalgebra-b.p005
Misalkan $\mathcal{A}$, $\mathcal{A}'$, dan seterusnya sembarang kategori. Pertama-tama, ambil funktor
% segment-id: o014.aljabr2.chapter9.coendomorphism-coalgebra-b.eq004
\[ \omega_i: \mathcal{A} \to \cated{Mod}B, \quad \omega'_i: \mathcal{A}' \to \cated{Mod}B, \quad i =1, 2. \]
Definisikan $\omega_i \boxtimes \omega'_i: \mathcal{A} \times \mathcal{A}' \to \cated{Mod}B$ dengan memetakan objek $(X, X')$ ke $\omega_i(X) \dotimes{B} \omega'_i(X')$.

% segment-id: o014.aljabr2.chapter9.coendomorphism-coalgebra-b.p006
Dengan mengandaikan bahwa $\coHom$ yang dibahas ada, terapkan sifat universal pada keluarga morfisme
% segment-id: o014.aljabr2.chapter9.coendomorphism-coalgebra-b.eq005
\begin{align*}
	\omega_2(X) \dotimes{B} \omega'_2(X') & \to \omega_1(X) \dotimes{B} \coHom_B(\omega_1, \omega_2) \dotimes{B} \omega'_1(X') \dotimes{B} \coHom_B(\omega'_1, \omega'_2) \\
	& \simeq \omega_1(X) \dotimes{B} \omega'_1(X') \dotimes{B} \coHom_B(\omega_1, \omega_2) \dotimes{B} \coHom_B(\omega'_1, \omega'_2)
\end{align*}
untuk memperoleh morfisme kanonik
% segment-id: o014.aljabr2.chapter9.coendomorphism-coalgebra-b.eq006
\[ \nu: \coHom_B(\omega_1 \boxtimes \omega'_1, \omega_2 \boxtimes \omega'_2) \to \coHom_B(\omega_1, \omega_2) \dotimes{B} \coHom_B(\omega'_1, \omega'_2), \]
yang membuat diagram berikut komutatif untuk setiap $(X, X') \in \Obj(\mathcal{A} \times \mathcal{A}')$:
% segment-id: o014.aljabr2.chapter9.coendomorphism-coalgebra-b.fig001
\[\begin{tikzpicture}
	\node (A) {$\omega_2(X) \dotimes{B} \omega'_2(X')$};
	\node (B) [right=of A] {$\omega_1(X) \dotimes{B} \omega'_1(X') \dotimes{B} \coHom_B(\omega_1 \boxtimes \omega'_1, \omega'_2 \boxtimes \omega'_2)$};
	\node (C) [below=of B] {$\omega_1(X) \dotimes{B} \omega'_1(X') \dotimes{B} \coHom_B(\omega_1, \omega_2) \dotimes{B} \coHom_B(\omega'_1, \omega'_2)$.};
	
	\draw[->] (A) -- (B) node[midway, above] {\scriptsize kanonik};
	\draw[->] (A) -- (C.north west) node[midway, below left] {\scriptsize kanonik};
	\draw[->] (B) -- (C) node[midway, auto] {\scriptsize $\identity \otimes \identity \otimes \nu$};
\end{tikzpicture}\]

% segment-id: o014.aljabr2.chapter9.coendomorphism-coalgebra-b.lem001
\begin{lemma}
	Dengan notasi di atas, andaikan pula bahwa $\omega_i$ dan $\omega'_i$ semuanya bernilai dalam $\catesubd{Mod}{\mathrm{pfg}}B$. Maka $\omega_i \boxtimes \omega'_i$ juga bernilai di dalamnya dan, dalam keadaan ini, $\nu$ merupakan isomorfisme. Mengikuti Konvensi \ref{con:coend-bracket}, isomorfisme tersebut dapat dideskripsikan secara konkret sebagai
	% segment-id: o014.aljabr2.chapter9.coendomorphism-coalgebra-b.eq007
	\begin{equation*}
		\nu: \left[ (\phi \otimes \phi') \otimes (x \otimes x') \right] \mapsto [\phi \otimes x] \otimes [\phi' \otimes x'].
	\end{equation*}
\end{lemma}
% segment-id: o014.aljabr2.chapter9.coendomorphism-coalgebra-b.proof001
\begin{proof}
	Karena $\catesubd{Mod}{\mathrm{pfg}}B$ tertutup terhadap $\otimes_B$, funktor $\omega \boxtimes \omega'$ bernilai di dalamnya. Dengan kembali ke konstruksi \eqref{eqn:cogebra-L} dan Proposisi \ref{prop:cogebra-L-universal}, deskripsi konkret untuk $\nu$ langsung diperoleh.
	
	Selanjutnya, kita tunjukkan bahwa $\nu$ merupakan isomorfisme. Menurut konstruksi dalam Proposisi \ref{prop:cogebra-L-universal}, $\coHom_B(\omega_1, \omega_2)$ dapat dinyatakan secara sesuai sebagai $\varinjlim_{X, Y} \omega_1(Y)^\vee \dotimes{B} \omega_2(X)$, dengan morfisme-morfisme terkait berasal dari \eqref{eqn:cogebra-L}; hal yang sama tentu berlaku untuk $\coHom_B$ lain yang terlibat dalam $\nu$. Oleh karena itu, pernyataan bahwa $\nu$ merupakan isomorfisme ekuivalen dengan pernyataan bahwa morfisme yang jelas
	% segment-id: o014.aljabr2.chapter9.coendomorphism-coalgebra-b.eq008
	\begin{multline*}
		\varinjlim_{X, X', Y, Y'} \omega_1(Y)^\vee \dotimes{B} \omega'_1(Y')^\vee \dotimes{B} \omega_2(X) \dotimes{B} \omega'_2(X') \\
		\to \varinjlim_{X, Y} \left( \omega_1(Y)^\vee \dotimes{B} \omega_2(X) \right) \dotimes{B} \varinjlim_{X', Y'} \left( \omega'_1(Y')^\vee \dotimes{B} \omega'_2(X') \right)
	\end{multline*}
	merupakan isomorfisme. Hal ini pada gilirannya direduksi ke fakta standar dalam teori modul bahwa $\dotimes{B}$ mempertahankan semua $\varinjlim$ kecil \cite[Proposisi 6.9.2]{Li1}.
\end{proof}

% segment-id: o014.aljabr2.chapter9.coendomorphism-coalgebra-b.p007
Kita lanjutkan pembahasan sebelumnya. Tinjau funktor-funktor
% segment-id: o014.aljabr2.chapter9.coendomorphism-coalgebra-b.eq009
\begin{gather*}
	\otimes: \mathcal{A} \times \mathcal{A}' \to \mathcal{A}'' , \\
	\omega_i: \mathcal{A} \to \catesubd{Mod}{\mathrm{pfg}}B, \quad \omega'_i: \mathcal{A}' \to \catesubd{Mod}{\mathrm{pfg}}B, \quad \omega''_i: \mathcal{A}'' \to \catesubd{Mod}{\mathrm{pfg}}B \quad i =1, 2,
\end{gather*}
beserta isomorfisme $\theta_i: \omega_i \boxtimes \omega'_i \rightiso \omega''_i \circ \otimes$, yang versi dualnya dinotasikan dengan $\theta_i^\vee$. Persiapan di atas memberikan homomorfisme $B$-linear
% segment-id: o014.aljabr2.chapter9.coendomorphism-coalgebra-b.eq010
\begin{equation}\label{eqn:L-tensor}
	\begin{tikzcd}[row sep=tiny, column sep=small]
		\coHom_B(\omega_1, \omega_2) \dotimes{B} \coHom_B(\omega'_1, \omega'_2) \arrow[r, "{\nu^{-1}}", "{\theta_1, \theta_2}"' inner sep=0.6em] & \coHom_B(\omega''_1 \circ \otimes, \; \omega''_2 \circ \otimes) \arrow[r, "{\text{Catatan \ref{rem:L-cogebra-functoriality}}}" inner sep=0.6em] & \coHom_B(\omega''_1, \omega''_2) \\
		{[\phi \otimes x]} \otimes {[\phi' \otimes x']} \arrow[mapsto, rr] & & {\left[ \theta_1^\vee(\phi \otimes \phi'), \theta_2(x \otimes x') \right]} .
	\end{tikzcd}
\end{equation}
% segment-id: o014.aljabr2.chapter9.coendomorphism-coalgebra-b.p008
Diagram komutatif untuk $\nu$ di atas menyiratkan bahwa diagram berikut komutatif; perinciannya diserahkan kepada pembaca.
% segment-id: o014.aljabr2.chapter9.coendomorphism-coalgebra-b.eq011
\begin{equation}\label{eqn:L-bialgebra-prod}
	\begin{tikzcd}[row sep=small, column sep=small]
		\omega''_2(X \otimes X') \arrow[d, "{\theta_2^{-1}}", "\sim"' sloped] \arrow[r, "\text{kanonik}"] & \omega''_1(X \otimes X') \dotimes{B} \coHom_B(\omega''_1, \omega''_2) \\
		\omega_2(X) \dotimes{B} \omega'_2(X') \arrow[d, "\text{kanonik}" inner sep=0.6em] &  \\
		\omega_1(X) \dotimes{B} \omega'_1(X') \dotimes{B} \coHom_B(\omega_1, \omega_2) \dotimes{B} \coHom_B(\omega'_1, \omega'_2) \arrow[r, "\text{\eqref{eqn:L-tensor}}"' inner sep=0.6em] & \omega_1(X) \dotimes{B} \omega'_1(X') \dotimes{B} \coHom_B(\omega''_1, \omega''_2). \arrow[uu, "{\theta_1 \otimes \identity}"', "\sim" sloped]
	\end{tikzcd}
\end{equation}

% segment-id: o014.aljabr2.chapter9.coendomorphism-coalgebra-b.p009
Sekarang andaikan lebih lanjut bahwa $\mathcal{A} = \mathcal{A}' = \mathcal{A}''$ merupakan kategori monoidal, $\otimes$ adalah perkaliannya, dan $\omega''_i = \omega'_i = \omega_i$ merupakan funktor monoidal ($i=1,2$). Tuliskan $\mathcal{A}_0$ untuk kategori yang hanya memiliki satu objek dan satu morfisme (yakni identitas), dan definisikan $\iota: \mathcal{A}_0 \to \mathcal{A}$ dengan memetakan objek tunggal tersebut ke $\munit$. Manipulasi langsung terhadap definisi memberikan
% segment-id: o014.aljabr2.chapter9.coendomorphism-coalgebra-b.eq012
\[ \coHom_B(\omega_1 \iota, \omega_2 \iota) \simeq B \dotimes{B} B \simeq B. \]
Funktorialitas pada Catatan \ref{rem:L-cogebra-functoriality} kemudian memberikan pemetaan kanonik
% segment-id: o014.aljabr2.chapter9.coendomorphism-coalgebra-b.eq013
\begin{equation}\label{eqn:L-bialgebra-unit}
	B \simeq \coHom_B(\omega_1 \iota, \omega_2 \iota) \to \coHom_B(\omega_1, \omega_2), \quad b \mapsto [b \otimes 1] = [1 \otimes b].
\end{equation}

% segment-id: o014.aljabr2.chapter9.coendomorphism-coalgebra-b.prop001
\begin{proposition}\label{prop:L-bialgebra}
	Dalam situasi di atas, andaikan $\mathcal{A} = \mathcal{A}' = \mathcal{A}''$ merupakan kategori monoidal, $\omega_1$ dan $\omega_2$ merupakan funktor monoidal, serta $\omega''_i = \omega'_i = \omega_i$.
	\begin{enumerate}[(i)]
% segment-id: o014.aljabr2.chapter9.coendomorphism-coalgebra-b.item001
		\item Homomorfisme yang bersesuaian $\coHom_B(\omega_1, \omega_2) \dotimes{B} \coHom_B(\omega_1, \omega_2) \to \coHom_B(\omega_1, \omega_2)$ menjadikan $\coHom_B(\omega_1, \omega_2)$ suatu aljabar-$B$, dengan \eqref{eqn:L-bialgebra-unit} sebagai unitnya.
% segment-id: o014.aljabr2.chapter9.coendomorphism-coalgebra-b.item002
		\item Dalam kasus khusus $\omega_1 = \omega = \omega_2$, struktur ini menjadikan $\coEnd_B(\omega)$ suatu bialjabar.
% segment-id: o014.aljabr2.chapter9.coendomorphism-coalgebra-b.item003
		\item Jika $\mathcal{A}$ adalah kategori monoidal simetris dan $\omega_1$, $\omega_2$ kompatibel dengan struktur kepang, maka $\coHom_B(\omega_1, \omega_2)$ merupakan aljabar-$B$ komutatif.
	\end{enumerate}
\end{proposition}
% segment-id: o014.aljabr2.chapter9.coendomorphism-coalgebra-b.proof002
\begin{proof}
	Karena perkalian pada $\coHom_B(\omega_1, \omega_2)$ memiliki karakterisasi alami berdasarkan \eqref{eqn:L-tensor} dan \eqref{eqn:L-bialgebra-prod}, seluruh pernyataan sisanya mencerminkan struktur monoidal dan hanya memerlukan argumen formal; perinciannya dihilangkan.
\end{proof}
